\chapter{Pemetaan Konformal dan Teorema Pemetaan Riemann}
\label{ch:b3:conformal}

Sebuah bijeksi \omterm{def:b3:holomorphic:holo}{holomorfik} antara dua daerah mengangkut seluruh
analisis kompleks dari yang satu ke yang lain: sehingga pemetaan semacam itu ---
yang \emph{konformal}, sebab ia mengawetkan sudut --- merupakan
isomorfisma dunia \omterm{def:b3:holomorphic:holo}{holomorfiknya}. Bab ini menggolongkannya
di tempat penggolongannya mungkin (yakni cakram dan bidangnya:
lema Schwarz adalah kuncinya, yaitu sebuah ketaksamaan yang dahsyat
kuasanya), membangun teori kekompakan keluarga \omterm{def:b3:holomorphic:holo}{holomorfiknya}
(lewat Montel), lalu membuktikan teorema keberadaan yang paling dalam pada
pokok bahasan ini: bahwa \emph{setiap} subdaerah sejati $\C$ yang \omterm{def:b3:conformal:simplyconnected}{terhubung
sederhana}, betapapun bergerigi batasnya, \omterm{def:b3:conformal:conformal}{setara secara konformal}
dengan cakram satuannya. Kita menutupnya dengan \omterm{prop:b3:conformal:harmonicholo}{fungsi harmonik} dan
\omterm{thm:b3:conformal:poisson}{kernel Poisson}, sambil menyelesaikan soal Dirichlet pada cakramnya ---
yakni hasil analitis geometri konformalnya. Sepanjang bab ini,
$\mathbb D = D(0,1)$ dan $\mathbb H = \{\operatorname{Im}z >
0\}$.

\section{Pemetaan konformal; transformasi Möbius}

\begin{definition}\label{def:b3:conformal:conformal}
Sebuah \emph{pemetaan konformal}\index{pemetaan konformal} (atau
biholomorfisma) antara \omterm{def:b3:topology:topology}{himpunan terbuka} adalah bijeksi \omterm{def:b3:holomorphic:holo}{holomorfik};
dan balikannya dengan sendirinya \omterm{def:b3:holomorphic:holo}{holomorfik}
(\cref{cor:b3:residues:openmapping}). Dua daerah disebut
\emph{setara secara konformal} jika pemetaan semacam itu ada; sedangkan
$\operatorname{Aut}(\Omega)$ menandai grup pemetaan diri
yang konformal. Di tempat $f' \neq 0$ --- yakni di mana-mana, bagi $f$ yang injektif
(lihat bukti \cref{cor:b3:residues:openmapping}) ---
diferensialnya adalah perkalian dengan $f'(z) \neq 0$: yakni sebuah
keserupaan, sehingga pemetaan konformal mengawetkan sudut antara kurva,
termasuk arahnya.
\end{definition}

\begin{example}[Transformasi Möbius]
\label{ex:b3:conformal:mobius}
Untuk $\bigl(\begin{smallmatrix}a & b\\ c &
d\end{smallmatrix}\bigr) \in GL_2(\C)$, \emph{transformasi
Möbius}\index{transformasi Möbius} $z \mapsto
\frac{az + b}{cz + d}$ bersifat konformal dari
$\C\setminus\{-d/c\}$ ke $\C\setminus\{a/c\}$ (dengan balikan
bertipe sama, dari matriks balikannya; sedangkan komposisinya
bersesuaian dengan hasil kali matriksnya). Sedangkan \emph{pemetaan
Cayley}\index{transformasi Cayley}
\[
\varphi(z) = \frac{z - \iu}{z + \iu}
\]
memetakan $\mathbb H$ secara konformal ke $\mathbb D$: sebab sungguh $\abs{z
- \iu} < \abs{z + \iu}$ tepat saat $z$ lebih dekat ke $\iu$
daripada ke $-\iu$, yakni $\operatorname{Im}z > 0$; dan balikannya adalah
$w \mapsto \iu\frac{1 + w}{1 - w}$. Pemetaan Möbius mengirim
keluarga lingkaran-dan-garis ke dirinya sendiri (\cref{exo:b3:conformal:1}).
\end{example}

\begin{example}[Pemetaan Joukowski]\label{ex:b3:conformal:joukowski}
Di luar Möbius, \omterm{def:b3:conformal:conformal}{pemetaan konformal} yang paling berguna pada
matematika terapan klasik adalah
\[
J(z) = \frac12\Bigl(z + \frac1z\Bigr) .
\]
Pada bagian luar $\Omega = \{\abs z > 1\}$ cakram satuannya,
$J$ bersifat injektif: sebab $J(z) = J(w)$ memberikan $(z - w)(1 -
\frac1{zw}) = 0$ padahal $\abs{zw} > 1$. Sedangkan turunannya $J'(z) =
\frac12(1 - z^{-2})$ lenyap hanya di $z = \pm1$, yakni pada
batasnya: jadi $J$ merupakan kesetaraan konformal dari $\Omega$ ke
petanya, yakni $\C\setminus\intcc{-1}1$ --- sedangkan lingkaran
satuannya sendiri terlipat dua-ke-satu ke ruasnya
(sebab $J(\eu^{\iu\theta}) = \cos\theta$). Jadi bagian luar sebuah
\emph{ruas}, yakni bidang teriris tanpa batas yang mulus, secara
konformal adalah bagian luar sebuah cakram: sehingga sudut bukanlah rintangan
bagi kesetaraan konformal, melainkan bagi kemulusan batasnya. Sedangkan peta
lingkaran yang melalui $\pm1$ tetapi tak berpusat di sana berupa kurva
berbentuk aerofoil, dan mengomposisikan $J$ dengan pemetaan Möbius mengangkut
aliran di sekitar sebuah silinder --- yang dapat dihitung dengan tangan --- ke aliran
di sekitar sebuah sayap: sehingga pada paruh pertama abad kedua puluh,
contoh inilah aerodinamikanya. Ia juga pintu menuju
Chebyshev: sebab $J$ mengonjugatkan $z \mapsto z^n$ ke polinomial
Chebyshev $T_n$ (\cref{pb:b3:hilbert:1}, Bagian V), karena
$J(z^n) = \cos(n\theta)$ bila $z = \eu^{\iu\theta}$.
\end{example}

\section{Lema Schwarz dan grup automorfismanya}

\begin{theorem}[Lema Schwarz]\label{thm:b3:conformal:schwarz}
Misalkan $f \colon \mathbb D \to \mathbb D$ \omterm{def:b3:holomorphic:holo}{holomorfik} dengan
$f(0) = 0$. Maka\index{lema Schwarz}
\[
\abs{f(z)} \leq \abs z \ \ (z \in \mathbb D)
\qquad\text{dan}\qquad \abs{f'(0)} \leq 1 ,
\]
dan jika $\abs{f(z_0)} = \abs{z_0}$ untuk satu $z_0 \neq 0$, atau
$\abs{f'(0)} = 1$, maka $f(z) = \eu^{\iu\theta}z$ merupakan sebuah
rotasi.
\end{theorem}

\begin{proof}
Fungsi $g(z) = f(z)/z$ diperluas secara \omterm{def:b3:holomorphic:holo}{holomorfik} ke $\mathbb D$ (sebab
kesingularannya di $0$ terhapuskan: karena $g$ terbatas di dekat $0$,
\cref{thm:b3:residues:riemanncw}; dan nilainya di $0$ adalah
$f'(0)$). Pada $\abs z = r < 1$: $\abs g \leq \frac1r$, sehingga menurut
asas maksimumnya (\cref{thm:b3:holomorphic:maximum})
$\abs g \leq \frac1r$ pada $\bar D(0,r)$; lalu biarkan $r \to 1$: $\abs
g \leq 1$ pada $\mathbb D$, yang merupakan kedua ketaksamaannya. Sedangkan kesamaan
di titik \omterm{def:b3:topology:interior}{interiornya} membuat $\abs g$ mencapai maksimum di
\omterm{def:b3:topology:interior}{interiornya}: sehingga $g$ konstan bermodulus $1$.
\end{proof}

\begin{theorem}[Automorfisma cakramnya]
\label{thm:b3:conformal:autdisc}
Untuk $a \in \mathbb D$, \emph{faktor
Blaschke}\index{faktor Blaschke}
\[
\varphi_a(z) = \frac{z - a}{1 - \bar a z}
\]
merupakan automorfisma $\mathbb D$ yang menukar $a$ dan $0$, dengan
$\varphi_a^{-1} = \varphi_{-a}$. Dan setiap automorfisma
$\mathbb D$ berbentuk $\eu^{\iu\theta}\varphi_a$ untuk $\theta
\in \R/2\pi\Z$ dan $a \in \mathbb D$ yang tunggal.
\end{theorem}

\begin{proof}
Pada $\abs z = 1$: $\abs{1 - \bar az} = \abs{\bar z}\abs{1 -
\bar az} = \abs{\bar z - \bar a\abs z^2} = \abs{\bar z - \bar
a} = \abs{z - a}$, sehingga $\abs{\varphi_a} = 1$ di sana; lalu menurut
asas maksimumnya $\varphi_a(\mathbb D) \subseteq \bar{\mathbb
D}$, dan keterbukaannya menaruh petanya di $\mathbb D$. Sedangkan
identitas aljabar $\varphi_{-a}\circ\varphi_a = \mathrm{id}$ (lewat
perhitungan langsung) menunjukkan kebijektifannya. Kini misalkan $f \in
\operatorname{Aut}(\mathbb D)$ dan $a = f^{-1}(0)$: maka $g =
f\circ\varphi_{-a}$ merupakan automorfisma yang menetapkan $0$. Lalu Schwarz
yang diterapkan pada $g$ dan pada $g^{-1}$: $\abs{g(z)} \leq \abs z$ dan
$\abs{g^{-1}(w)} \leq \abs w$, sehingga $\abs{g(z)} = \abs z$: jadi
rotasi, $g = \eu^{\iu\theta}\,\mathrm{id}$, yakni $f =
\eu^{\iu\theta}\varphi_a$. Untuk ketunggalannya: $a = f^{-1}(0)$ dan
$\theta$ dari penilaian bertipe $f'$ (atau dari $f(0) =
-\eu^{\iu\theta}a$ dan satu nilai lagi).
\end{proof}

\begin{theorem}[Automorfisma bidangnya]
\label{thm:b3:conformal:autplane}
$\operatorname{Aut}(\C) = \{z \mapsto az + b : a \in \C^*,\ b
\in \C\}$. Karena itu $\C$ dan $\mathbb D$ tak
\omterm{def:b3:conformal:conformal}{setara secara konformal}.
\end{theorem}

\begin{proof}
Misalkan $f \in \operatorname{Aut}(\C)$ lalu tinjaulah $g(z) =
f(1/z)$ pada $\C^*$: yakni \omterm{def:b3:holomorphic:holo}{fungsi holomorfik} yang berkesingularan
terpencil di $0$. Seandainya ia hakiki, maka
Casorati--Weierstrass (\cref{thm:b3:residues:riemanncw}) akan
membuat $g\bigl(D(0,\varepsilon)\setminus\{0\}\bigr)$ padat,
sedangkan $f(D(0, 1))$ terbuka dan saling lepas dengannya (sebab $f$
injektif: karena kedua himpunannya peta himpunan yang saling lepas) ---
yang mustahil bagi sebuah \omterm{def:b3:topology:interior}{himpunan padat} dan sebuah \omterm{def:b3:topology:topology}{himpunan terbuka} tak kosong. Jadi $0$
berupa kutub atau terhapuskan bagi $g$, yakni $\abs{f(z)}$ bertumbuh paling banyak
polinomial: sehingga $f$ merupakan polinomial
(\cref{exo:b3:holomorphic:4}(a)). Lalu keinjektifannya memaksa derajat
$1$: sebab polinomial berderajat lebih tinggi entah berakar rangkap pada
$f - c$ di suatu tempat (sehingga $f'$ lenyap) entah berprapeta beberapa
yang berbeda (menurut d'Alembert--Gauss, \cref{pb:b3:holomorphic:1});
dan bagaimanapun keinjektifannya gagal. Akhirnya, sebuah \omterm{def:b3:conformal:conformal}{pemetaan konformal} $\C \to
\mathbb D$ akan berupa fungsi utuh yang terbatas: jadi konstan
(menurut Liouville) --- sehingga tak ada kesetaraannya.
\end{proof}

\section{Teorema Montel}

\begin{theorem}[Montel]\label{thm:b3:conformal:montel}
Misalkan $\mathcal F \subseteq \mathcal H(\Omega)$ \emph{terbatas
lokal}: yakni setiap titiknya berpersekitaran yang $\sup_{f\in
\mathcal F}\sup\abs f < \infty$ padanya. Maka setiap barisan
$\mathcal F$ mempunyai subbarisan yang konvergen seragam pada setiap
himpunan bagian \omterm{def:b3:topology:compact}{kompak} $\Omega$ (dengan limit yang
\omterm{def:b3:holomorphic:holo}{holomorfik}).\index{teorema Montel}
\end{theorem}

\begin{proof}
Untuk \omterm{def:b3:topology:continuity}{kekontinuan} setara lokalnya: jika $\abs f \leq M$ pada $D(a, 2r)
\subseteq \Omega$ untuk setiap $f \in \mathcal F$, maka rumus
Cauchy memberikan, untuk $z, z' \in D(a, r)$,
\[
\abs{f(z) - f(z')} = \frac{\abs{z - z'}}{2\pi}
\Bigl|\int_{C_{2r}}\frac{f(w)\,\dd w}{(w-z)(w-z')}\Bigr|
\leq \frac{\abs{z - z'}\,2\pi\cdot2r\,M}{2\pi\,r^2}
= \frac{2M}{r}\,\abs{z - z'} :
\]
yakni sebuah batas Lipschitz yang seragam. Lalu habiskan $\Omega$ oleh \omterm{def:b3:topology:compact}{kompak} $K_m$;
sebab setiap $K_m$ terselimuti berhingga banyak cakram semacam itu, sehingga
$\mathcal F$ terbatas seragam dan berkekontinuan setara pada
$K_m$: sehingga Arzelà--Ascoli (\cref{thm:b3:complete:ascoli})
memetik subbarisan yang konvergen seragam pada $K_m$;
lalu diagonalkan atas $m$. Sedangkan limitnya \omterm{def:b3:holomorphic:holo}{holomorfik} menurut
\cref{thm:b3:holomorphic:weierstrassconv}.
\end{proof}

\section{Teorema pemetaan Riemann}

\begin{definition}\label{def:b3:conformal:simplyconnected}
Sebuah $\Omega \subseteq \C$ yang terbuka \omterm{def:b3:topology:connected}{terhubung} disebut \emph{terhubung
sederhana}\index{daerah terhubung sederhana} (dalam arti homologinya,
yang mencukupi bagi seluruh keperluan kita) jika
$\operatorname{Ind}_\gamma(w) = 0$ untuk setiap siklus $\gamma$ di
$\Omega$ dan setiap $w \notin \Omega$ --- yakni ``tak ada siklus
$\Omega$ yang mengelilingi lubang''. Lalu menurut teorema Cauchy globalnya
(\cref{thm:b3:residues:globalcauchy}) dan
\cref{prop:b3:holomorphic:primitive}, pada $\Omega$ semacam itu
\emph{setiap \omterm{def:b3:holomorphic:holo}{fungsi holomorfik} berantiturunan}; sehingga setiap
$f \in \mathcal H(\Omega)$ yang bebas akar mempunyai logaritma
\omterm{def:b3:holomorphic:holo}{holomorfik} (yakni $\exp\circ$(antiturunan $f'/f$), yang disesuaikan oleh sebuah
konstanta, sebab $(f\eu^{-L})' = 0$) beserta akar ke-$n$ \omterm{def:b3:holomorphic:holo}{holomorfiknya}
$\eu^{L/n}$.\index{logaritma holomorfik}
\end{definition}

\begin{theorem}[Teorema pemetaan Riemann]
\label{thm:b3:conformal:rmt}
Setiap $\Omega \subsetneq \C$ yang terbuka \omterm{def:b3:conformal:simplyconnected}{terhubung sederhana} dengan $\Omega
\neq \varnothing$ \omterm{def:b3:conformal:conformal}{setara secara konformal} dengan $\mathbb D$;
dan diberikan $z_0 \in \Omega$, ada tepat satu $f \colon
\Omega \to \mathbb D$ yang konformal dengan $f(z_0) = 0$ dan $f'(z_0) >
0$.\index{teorema pemetaan Riemann}
\end{theorem}

\begin{proof}
\emph{Langkah 0: keluarganya tak kosong.} Pilihlah $b \notin \Omega$:
maka $z - b$ bebas akar pada $\Omega$, sehingga ia berakar kuadrat
\omterm{def:b3:holomorphic:holo}{holomorfik} $h$ (dengan $h^2 = z - b$). Lalu $h$ injektif (sebab $h(z) =
h(z')$ berkuadrat $z = z'$), dan jika $w \in h(\Omega)$ maka $-w
\notin h(\Omega)$ (sebab $h(z) = -h(z')$ juga berkuadrat $z = z'$,
yang memberikan $w = -w = 0$, yang mustahil karena $h$ bebas akar). Lalu karena
$h(\Omega)$ terbuka, ia memuat sebuah cakram $D(h(z_0), \rho)$;
sehingga $D(-h(z_0), \rho) \cap h(\Omega) = \varnothing$, yakni
$\abs{h(z) + h(z_0)} \geq \rho$ untuk setiap $z \in \Omega$.
Karena itu
\[
g(z) = \frac{\rho}{2\,\bigl(h(z) + h(z_0)\bigr)}
\]
bersifat \omterm{def:b3:holomorphic:holo}{holomorfik}, injektif (yakni pemetaan Möbius yang dikomposisikan dengan
$h$ yang injektif), dengan $\abs g \leq \frac12 < 1$. Lalu dengan mengomposisikannya dengan sebuah \omterm{thm:b3:conformal:autdisc}{faktor Blaschke}
(\cref{thm:b3:conformal:autdisc}) untuk memindahkan $g(z_0)$ ke $0$,
keluarganya
\[
\mathcal F = \{f \colon \Omega \to \mathbb D \text{
holomorfik, injektif, } f(z_0) = 0\}
\]
bersifat tak kosong.

\emph{Langkah 1: sebuah unsur ekstremal.} Misalkan $s =
\sup_{\mathcal F}\abs{f'(z_0)} \in \intoc0{+\infty}$ (yang $> 0$: sebab
anggotanya injektif, sehingga $f'(z_0) \neq 0$). Ambil $f_n \in
\mathcal F$ dengan $\abs{f_n'(z_0)} \to s$: keluarganya
terbatas oleh $1$, sehingga Montel (\cref{thm:b3:conformal:montel})
memetik $f_n \to f$ secara seragam pada \omterm{def:b3:topology:compact}{kompak}; lalu $f$ bersifat
\omterm{def:b3:holomorphic:holo}{holomorfik} dengan $f(z_0) = 0$ dan $\abs{f'(z_0)} = s$
(lewat \cref{thm:b3:holomorphic:weierstrassconv} untuk
turunannya), sehingga khususnya $f$ tak konstan; lalu $f$ bersifat
injektif menurut Hurwitz (\cref{exo:b3:residues:8}(b)), dan
$f(\Omega) \subseteq \bar{\mathbb D}$, jadi $\subseteq
\mathbb D$ (lewat pemetaan terbukanya). Jadi $f \in \mathcal F$ mencapai
supremumnya: sehingga $s < \infty$.

\emph{Langkah 2: pemetaan ekstremalnya bersifat pada.} Andaikan $a \in
\mathbb D\setminus f(\Omega)$. Maka angkutan Blaschke
$\varphi_a\circ f$ bebas akar pada $\Omega$ yang \omterm{def:b3:conformal:simplyconnected}{terhubung
sederhana}: sehingga ia berakar kuadrat \omterm{def:b3:holomorphic:holo}{holomorfik} $F$ (dengan
$F(\Omega) \subseteq \mathbb D$, sebab $\abs F^2 =
\abs{\varphi_a\circ f} < 1$), yang injektif (sebab kuadratnya
membedakan). Lalu normalkan: $G = \varphi_{F(z_0)}\circ F \in
\mathcal F$. Dengan membalikkannya: $f = \varphi_{-a}\circ s_2 \circ
\varphi_{-F(z_0)}\circ G$ dengan $s_2(w) = w^2$; sedangkan pemetaan $\Psi
= \varphi_{-a}\circ s_2\circ\varphi_{-F(z_0)} \colon \mathbb D
\to \mathbb D$ bersifat \omterm{def:b3:holomorphic:holo}{holomorfik} dengan $\Psi(0) = f(z_0) = 0$ dan
\emph{bukan} sebuah rotasi (sebab ia tak injektif: karena $s_2$ tidak).
Lalu lema Schwarz (pada kasus sejatinya): $\abs{\Psi'(0)} < 1$, dan
aturan rantai $f = \Psi\circ G$ memberikan $\abs{f'(z_0)} =
\abs{\Psi'(0)}\,\abs{G'(z_0)} < \abs{G'(z_0)}$ ---
yang bertentangan dengan kemaksimalannya (perhatikan $G \in \mathcal F$). Karena itu $f$
bersifat pada: yakni sebuah kesetaraan konformal.

\emph{Langkah 3: penormalan dan ketunggalannya.} Kalikan $f$ dengan
$\eu^{-\iu\arg f'(z_0)}$ untuk membuat $f'(z_0) > 0$ (dan ini tetap di
$\mathcal F$). Lalu jika $f_1, f_2$ sama-sama berhasil, maka $\psi =
f_2\circ f_1^{-1} \in \operatorname{Aut}(\mathbb D)$ menetapkan
$0$ dengan $\psi'(0) = f_2'(z_0)/f_1'(z_0) > 0$; sehingga menurut
\cref{thm:b3:conformal:autdisc} $\psi$ merupakan rotasi
$\eu^{\iu\theta}$ dengan $\eu^{\iu\theta} > 0$: jadi $\psi =
\mathrm{id}$.
\end{proof}

\begin{remark}\label{rem:b3:conformal:rmtscope}
Teoremanya merupakan pernyataan keberadaan murni yang jangkauannya
menakjubkan: sebuah bujur sangkar, sebuah setengah bidang, komplemen sebuah irisan,
daerah antara dua lingkaran yang bersinggungan, sebuah daerah berbatas fraktal
--- semuanya identik secara konformal dengan $\mathbb D$. Sedangkan yang tak
diberikannya: rumus apa pun (karena pemetaan eksplisit merupakan kekecualian:
\cref{exo:b3:conformal:5}), perilaku batasnya (yang ditangani teori yang lebih dalam
--- yakni teorema Carathéodory), maupun ketunggalan
perluasannya ke $\C$ atau ke daerah \omterm{def:b3:topology:connected}{terhubung} ganda: sebab
anulus $\{1 < \abs z < 2\}$ \emph{tak} secara konformal berupa
cakram tertusuk, dan anulus yang nisbah jari-jarinya berbeda
tak setara (yakni fakta yang sungguh lebih sukar).
\end{remark}

\section{Fungsi harmonik dan kernel Poisson}

\begin{proposition}\label{prop:b3:conformal:harmonicholo}
Misalkan $\Omega$ \omterm{def:b3:conformal:simplyconnected}{terhubung sederhana} dan $u \colon \Omega \to \R$
harmonik\index{fungsi harmonik} (yakni $\mathcal C^2$ dengan
$\Delta u = u_{xx} + u_{yy} = 0$). Maka $u =
\operatorname{Re}F$ untuk sebuah $F$ yang \omterm{def:b3:holomorphic:holo}{holomorfik} dan tunggal sampai sebuah
konstanta imajiner. Karena itu $u$ bersifat $\mathcal C^\infty$,
memenuhi sifat nilai rata-rata, dan menuruti asas
maksimumnya (yakni tak ada ekstremum \omterm{def:b3:topology:interior}{interior} sejati kecuali ia konstan).
\end{proposition}

\begin{proof}
Fungsi $g = u_x - \iu u_y$ memenuhi \omterm{prop:b3:holomorphic:cauchyriemann}{persamaan Cauchy--Riemann}
(dengan $P = u_x$ dan $Q = -u_y$: sebab $P_x = u_{xx} = -u_{yy} = Q_y$ dan
$P_y = u_{xy} = u_{yx} = -Q_x$) beserta turunan parsial yang \omterm{def:b3:topology:continuity}{kontinu}: sehingga $g$
bersifat \omterm{def:b3:holomorphic:holo}{holomorfik} (\cref{prop:b3:holomorphic:cauchyriemann}; sebab
keterturunan-$\R$-nya menyusul dari $\mathcal C^1$). Lalu misalkan $F_0$
sebuah antiturunan (lewat keterhubungan sederhananya,
\cref{def:b3:conformal:simplyconnected}); maka
$\operatorname{Re}F_0$ bergradien $(u_x, u_y)$
(sebab $F_0' = g$ terurai tepat menjadi itu lewat Cauchy--Riemann bagi
$F_0$), sehingga $u - \operatorname{Re}F_0$ konstan (sebab $\Omega$
\omterm{def:b3:topology:connected}{terhubung}): lalu sesuaikan $F = F_0 + c$. Sedangkan sifatnya berpindah
dari \cref{thm:b3:holomorphic:maximum} dan
\cref{exo:b3:holomorphic:10} (dan untuk asas maksimum yang
diterapkan pada $u$ sendiri, pakailah $\eu^{F}$ seperti di sana).
\end{proof}

\begin{theorem}[Rumus Poisson; soal Dirichlet pada
cakramnya]\label{thm:b3:conformal:poisson}
Untuk $0 \leq r < 1$ definisikan \emph{kernel
Poisson}\index{kernel Poisson}
\[
P_r(\theta) = \sum_{n\in\Z}r^{\abs n}\eu^{\iu n\theta}
= \frac{1 - r^2}{1 - 2r\cos\theta + r^2} \;>\; 0 .
\]
Lalu misalkan $g \colon \partial\mathbb D \to \R$ \omterm{def:b3:topology:continuity}{kontinu}, dan
tetapkan, untuk $z = r\eu^{\iu\varphi} \in \mathbb D$,
\[
u(z) = \frac1{2\pi}\int_0^{2\pi}
P_r(\varphi - t)\,g(\eu^{\iu t})\,\dd t .
\]
Maka $u$ bersifat harmonik pada $\mathbb D$ dan diperluas secara \omterm{def:b3:topology:continuity}{kontinu}
ke $\bar{\mathbb D}$ dengan nilai batas $g$: yakni satu-satunya
\omterm{prop:b3:conformal:harmonicholo}{fungsi harmonik} yang demikian.\index{soal Dirichlet}
\end{theorem}

\begin{proof}
\emph{Identitas kernelnya}: dengan menjumlahkan dua deret geometri,
\[
\sum_{n\in\Z}r^{\abs n}\eu^{\iu n\theta}
= \operatorname{Re}\frac{1 + r\eu^{\iu\theta}}{1 -
r\eu^{\iu\theta}}
= \frac{1 - r^2}{\abs{1 - r\eu^{\iu\theta}}^2},
\]
yang merupakan hasil bagi yang ditampilkan itu; sedangkan kepositifannya jelas, dan
$\frac1{2\pi}\int_0^{2\pi}P_r = 1$ (sebab hanya $n = 0$ yang bertahan).

\emph{Keharmonikannya}: dengan $z = r\eu^{\iu\varphi}$,
\[
u(z) = \operatorname{Re}\biggl[\frac1{2\pi}\int_0^{2\pi}
\frac{\eu^{\iu t} + z}{\eu^{\iu t} - z}\,
g(\eu^{\iu t})\,\dd t\biggr],
\]
(sebab kernel di dalam kurungnya berbagian real $P_r(\varphi - t)$:
hitunglah), dan kurungnya bersifat \omterm{def:b3:holomorphic:holo}{holomorfik} terhadap $z$ pada $\mathbb D$
(\cref{exo:b3:holomorphic:7}): sehingga $u$ merupakan bagian real sebuah
\omterm{def:b3:holomorphic:holo}{fungsi holomorfik}, jadi harmonik.

\emph{Nilai batasnya}: $P_r(\cdot)$ merupakan identitas
hampiran saat $r \to 1^-$: sebab bermassa $1$, dan untuk $\delta \leq
\abs\theta \leq \pi$, $P_r(\theta) \leq \frac{1 - r^2}{1 -
2r\cos\delta + r^2} \to 0$ secara seragam. Lalu pemisahan bakunya
(lewat \omterm{def:b3:topology:continuity}{kekontinuan} $g$ di dekat $\eu^{\iu\varphi_0}$ dan keterbatasannya
di tempat lain) memberikan $u(r\eu^{\iu\varphi}) \to
g(\eu^{\iu\varphi_0})$ saat $r\eu^{\iu\varphi} \to
\eu^{\iu\varphi_0}$, secara seragam terhadap titik batasnya: sehingga
perluasannya \omterm{def:b3:topology:continuity}{kontinu}. \emph{Untuk ketunggalannya}: selisih
dua solusinya bersifat harmonik pada $\mathbb D$, \omterm{def:b3:topology:continuity}{kontinu} pada
\omterm{def:b3:topology:interior}{penutupnya}, dan nol pada batasnya: sehingga menurut asas maksimumnya
(yang diterapkan pada $\pm$ selisihnya), ia lenyap.
\end{proof}

\begin{omfigure}
\begin{tikzpicture}[scale=1.05]
  \begin{scope}
    \clip (-2.1,-2.1) rectangle (2.1,2.1);
    \foreach \c in {0.2, 0.5, 0.8, 1.1, 1.4}{
      \draw[omDef!60] plot[domain=-2.05:2.05, samples=80]
        (\x, {\c/(2*\x)});
      \draw[omDef!60] plot[domain=-2.05:2.05, samples=80]
        (\x, {-\c/(2*\x)});
      \draw[omThm!60, samples=80, domain=\c:2.9]
        plot ({sqrt((\x*\x-\c*\c)/1)/1}, {0}) ;
    }
    % hyperbolas x^2 - y^2 = c
    \foreach \c in {0.4, 1, 1.6}{
      \draw[omThm!70] plot[domain=-1.32:1.32, samples=60]
        ({cosh(\x)*sqrt(\c)}, {sinh(\x)*sqrt(\c)});
      \draw[omThm!70] plot[domain=-1.32:1.32, samples=60]
        ({-cosh(\x)*sqrt(\c)}, {sinh(\x)*sqrt(\c)});
    }
  \end{scope}
  \draw[->] (-2.2,0) -- (2.3,0);
  \draw[->] (0,-2.2) -- (0,2.3);
  \node[font=\small, omThm] at (1.75,-1.9)
    {$\operatorname{Re}z^2 = c$};
  \node[font=\small, omDef] at (-1.35,1.95)
    {$\operatorname{Im}z^2 = c$};
\end{tikzpicture}

{\small Kurva aras $\operatorname{Re}z^2 = x^2 - y^2$
(yakni hiperbola merah) dan $\operatorname{Im}z^2 = 2xy$ (yakni hiperbola
biru) berpotongan tegak lurus jauh dari $0$: sehingga sebuah
\omterm{def:b3:conformal:conformal}{pemetaan konformal} ($z \mapsto z^2$, di tempat $z \neq 0$) mengawetkan
keortogonalan garis koordinatnya. Sedangkan di $z = 0$, tempat
turunannya lenyap, sudutnya justru berlipat dua.}
\end{omfigure}

\section{Latihan}

\begin{exercise}[$\star$]\label{exo:b3:conformal:1}
(a) Periksalah bahwa \omterm{ex:b3:conformal:mobius}{pemetaan Cayley} $\varphi(z) = \frac{z -
\iu}{z + \iu}$ merupakan bijeksi $\mathbb H \to \mathbb D$ dengan
balikan yang dinyatakan itu, lalu hitunglah peta $\iu$, $0$,
$1$, dan $\infty$ (lewat limitnya).
(b) Tunjukkan bahwa $z \mapsto 1/z$ memetakan lingkaran dan garis ke
lingkaran dan garis. \emph{(Tulislah persamaan bersamanya
$\alpha\abs z^2 + \bar\beta z + \beta\bar z + \gamma = 0$
dengan $\alpha, \gamma \in \R$.)} Lalu turunkan hal yang sama bagi setiap pemetaan
Möbius.
\end{exercise}

\begin{exercise}[$\star$]\label{exo:b3:conformal:2}
Misalkan $f \colon \mathbb D \to \mathbb D$ \omterm{def:b3:holomorphic:holo}{holomorfik}.
(a) Jika $f(0) = 0$ dan $f(a) = a$ untuk suatu $a \neq 0$, tunjukkan bahwa $f
= \mathrm{id}$.
(b) Jika $f$ merupakan automorfisma dengan dua titik tetap yang berbeda
di $\mathbb D$, tunjukkan bahwa $f = \mathrm{id}$
\emph{(konjugatkanlah dengan sebuah \omterm{thm:b3:conformal:autdisc}{faktor Blaschke} untuk menyusutkannya ke (a))}.
\end{exercise}

\begin{exercise}[$\star\star$]\label{exo:b3:conformal:3}
(Schwarz--Pick) Untuk $f\colon \mathbb D \to
\mathbb D$ yang \omterm{def:b3:holomorphic:holo}{holomorfik}, buktikan bahwa
\[
\frac{\abs{f'(z)}}{1 - \abs{f(z)}^2} \;\leq\;
\frac{1}{1 - \abs z^2}
\qquad (z \in \mathbb D),
\]
dengan kesamaannya (di satu titik, sehingga di mana-mana) jika dan hanya jika $f \in
\operatorname{Aut}(\mathbb D)$.
\emph{(Terapkan Schwarz pada $\varphi_{f(z)}\circ
f\circ\varphi_{-z}$.)} Tafsirannya: bahwa pemetaan diri yang \omterm{def:b3:holomorphic:holo}{holomorfik}
mengerutkan metrik hiperboliknya.
\end{exercise}

\begin{exercise}[$\star\star$]\label{exo:b3:conformal:4}
Carilah kesetaraan konformal yang eksplisit:
(a) dari jalur $\{0 < \operatorname{Im}z < \pi\} \to \mathbb H$;
(b) dari kuadran $\{\operatorname{Re}z > 0,
\operatorname{Im}z > 0\} \to \mathbb H$;
(c) dari setengah cakram $\mathbb D\cap\mathbb H \to$ sebuah kuadran,
lalu $\to \mathbb H$;
(d) dari $\mathbb D \to \mathbb D$ yang mengirim $\frac12$ ke $0$ dengan
turunan positif di sana.
\end{exercise}

\begin{exercise}[$\star\star$]\label{exo:b3:conformal:5}
(a) Tunjukkan bahwa tak ada \omterm{def:b3:conformal:conformal}{pemetaan konformal} $\C \to \mathbb D$ atau $\C \to
\mathbb H$, dan tak ada pula $\mathbb D \to \C$.
(b) Manakah di antara berikut yang \omterm{def:b3:conformal:conformal}{setara secara konformal} dengan
$\mathbb D$? Benarkanlah lewat \cref{thm:b3:conformal:rmt} atau lewat sebuah
rintangan: sebuah bujur sangkar; $\C\setminus\intoc{-\infty}0$;
$\mathbb D\setminus\{0\}$; dan $\{1 < \abs z < 2\}$.
\emph{(Untuk kedua yang terakhir: sebuah peta konformal cakram
tertusuknya akan diperluas melintasi tusukannya menurut
\cref{thm:b3:residues:riemanncw}(1) --- kembangkanlah ini.)}
\end{exercise}

\begin{exercise}[$\star\star$]\label{exo:b3:conformal:6}
Misalkan $\mathcal F = \{f \in \mathcal H(\mathbb D) : f(0) = 1,\
\operatorname{Re}f > 0\}$.
(a) Tunjukkan bahwa $\mathcal F$ terbatas lokal.
\emph{(Komposisikanlah dengan pemetaan bertipe Cayley $w \mapsto \frac{w -
1}{w + 1}$ yang mengirim setengah bidang kanannya ke $\mathbb D$, lalu
terapkan Schwarz.)}
(b) Turunkan \emph{batas Herglotz}: $\abs{f(z)} \leq
\frac{1 + \abs z}{1 - \abs z}$ untuk $f \in \mathcal F$, beserta
kemungkinan kesamaannya.
\end{exercise}

\begin{exercise}[$\star\star\star$]\label{exo:b3:conformal:7}
Di manakah bukti \cref{thm:b3:conformal:rmt} memakai masing-masing
hipotesisnya? Lacaklah: (i) keterhubungan sederhananya (dua kali); (ii)
$\Omega \neq \C$; (iii) keterhubungannya. Lalu tunjukkan bahwa
teoremanya \emph{gagal} untuk $\Omega = \C$ dan untuk anulusnya,
sambil menunjuk langkah mana pada buktinya yang patah pada masing-masing kasusnya.
\end{exercise}

\begin{exercise}[$\star\star$]\label{exo:b3:conformal:8}
Selesaikan soal Dirichlet pada $\mathbb D$ untuk data
batas: (a) $g(\eu^{\iu t}) = \cos t$; (b) $g(\eu^{\iu t}) =
\cos^2 t$; (c) $g = \mathbf 1_{\text{setengah lingkaran atas}}$ ---
dan untuk (c) hitunglah $u(0)$ lalu tafsirkan lewat sifat nilai
rata-ratanya. \emph{(Uraikan $g$ menjadi deret Fourier lalu pakailah
deret $P_r$: $u(r\eu^{\iu\varphi}) = \sum_n c_n(g)
r^{\abs n}\eu^{\iu n\varphi}$.)}
\end{exercise}

\begin{exercise}[$\star\star\star$]\label{exo:b3:conformal:9}
(Harnack) Misalkan $u \geq 0$ harmonik pada $\mathbb D$. Buktikan,
untuk $\abs z = r < 1$:
\[
\frac{1 - r}{1 + r}\,u(0) \;\leq\; u(z) \;\leq\;
\frac{1 + r}{1 - r}\,u(0)
\]
\emph{(batasilah \omterm{thm:b3:conformal:poisson}{kernel Poissonnya} antara $\frac{1-r}{1+r}$ dan
$\frac{1+r}{1-r}$; lalu terapkan penyajiannya pada cakram yang sedikit
lebih kecil lalu ambil limitnya)}. Lalu turunkan: bahwa \omterm{prop:b3:conformal:harmonicholo}{fungsi harmonik}
pada $\C$ yang terbatas dari bawah bersifat konstan.
\end{exercise}

\begin{exercise}[$\star\star$]\label{exo:b3:conformal:10}
Dengan memakai keinvarianan konformal keharmonikannya (yakni $u\circ f$ bersifat
harmonik bila $u$ harmonik dan $f$ \omterm{def:b3:holomorphic:holo}{holomorfik} --- buktikanlah
lewat \cref{prop:b3:conformal:harmonicholo} secara lokal), selesaikanlah
soal Dirichlet pada setengah bidang atasnya dengan data batas
$\mathbf 1_{\intoo{-\infty}0}$: tunjukkan bahwa
\[
u(x + \iu y) = \frac1\pi\,\arg(x + \iu y)
\qquad (\arg \in \intoo0\pi \text{ pada } \mathbb H)
\]
bersifat harmonik pada $\mathbb H$ (yakni bagian imajiner sebuah \omterm{def:b3:conformal:simplyconnected}{logaritma
holomorfik}) dengan limit batas yang dituntut di setiap $x \neq
0$, lalu angkutlah ia ke cakramnya lewat Cayley untuk menurunkan kembali
\cref{exo:b3:conformal:8}(c).
\end{exercise}

\begin{exercise}[$\star\star$]\label{exo:b3:conformal:11}
(Titik tetap dan iterasi di cakramnya) Misalkan $f\colon\mathbb
D \to \mathbb D$ \omterm{def:b3:holomorphic:holo}{holomorfik}.
(a) Tunjukkan bahwa jika $f$ mempunyai \emph{dua} titik tetap yang berbeda,
maka $f = \mathrm{id}$ \emph{(pindahkanlah salah satunya ke $0$ lewat sebuah
automorfisma lalu terapkan kasus kesamaan Schwarz)}.
(b) Andaikan $f(0) = 0$ dan $f$ bukan rotasi. Tunjukkan bahwa
iterasinya $f^{\circ n} \to 0$ secara seragam pada setiap \omterm{def:b3:topology:compact}{kompak}
$\bar D(0, r)$ dengan $r < 1$ \emph{(sebab Schwarz memberikan $\abs{f(z)}
\leq c_r\abs z$ pada $\bar D(0,r)$ dengan $c_r < 1$ --- benarkanlah
konstanta sejati ini lewat asas maksimum yang diterapkan pada
$f(z)/z$)}.
(c) Peragakan dengan $f(z) = \frac{z^2 + z}2$: yakni titik tetapnya,
dan laju kekonvergenan \omterm{def:b3:groups:action}{orbit} $z_0 =
\frac12$.
\end{exercise}

\begin{exercise}[$\star\star$]\label{exo:b3:conformal:12}
(Sekawan harmonik, secara konkret) Misalkan $u(x, y) = x^3 - 3xy^2
+ 2y$.
(a) Periksalah bahwa $u$ harmonik pada $\R^2$, lalu carilah semua
sekawan harmoniknya $v$ (yakni yang $u + \iu v$-nya \omterm{def:b3:holomorphic:holo}{holomorfik}) dengan
mengintegralkan persamaan Cauchy--Riemannya; lalu kenali $f(z) =
u + \iu v$ sebagai polinomial dalam $z$.
(b) Tunjukkan bahwa pada \omterm{def:b3:topology:topology}{himpunan terbuka} \emph{berbentuk bintang}, setiap
\omterm{prop:b3:conformal:harmonicholo}{fungsi harmonik} memiliki sekawan harmonik, yang tunggal sampai
sebuah konstanta penjumlah \emph{(sebab bentuk-$1$ $-u_y\,\dd x +
u_x\,\dd y$ bersifat tertutup; lalu pakailah mesin antiturunan
\cref{thm:b3:holomorphic:cauchy}, atau lema Poincaré
\cref{ch:b3:forms})}.
(c) Berikan contoh penyangkal bakunya pada
$\C^*$: bahwa $u = \ln\abs z$ tak bersekawan global --- lalu kaitkan
dengan bentuk sudutnya dan bilangan lilitannya
(\cref{ch:b3:forms}).
\end{exercise}

\section{Soal: teorema luas dan teorema seperempat
Koebe}

\begin{problem}[{Soal akhir pekan --- seberapa banyak yang harus
diselimuti sebuah pemetaan univalen?}]\label{pb:b3:conformal:1}
Sebuah fungsi \emph{univalen} adalah injeksi \omterm{def:b3:holomorphic:holo}{holomorfik}. Sedangkan
fungsi univalen yang ternormalkan pada cakramnya,
\[
\mathcal S = \bigl\{f \in \mathcal H(\mathbb D) \text{
injektif},\ f(z) = z + a_2z^2 + a_3z^3 + \cdots\bigr\},
\]
terkendala secara kaku: kita membuktikan ketaksamaan Bieberbach
$\abs{a_2} \leq 2$ lalu menurunkan \emph{teorema seperempat
Koebe}: bahwa peta setiap $f \in \mathcal S$ memuat
cakram $D(0, \frac14)$ --- yakni konstanta semesta yang tajam pada
geometri konformalnya.\index{teorema seperempat Koebe}\index{teorema
luas}

\medskip
\noindent\textbf{Bagian I --- Teorema luasnya.} Misalkan $g(w) = w
+ b_0 + \frac{b_1}w + \frac{b_2}{w^2} + \cdots$ \omterm{def:b3:holomorphic:holo}{holomorfik}
dan \emph{injektif} pada $\{\abs w > 1\}$.
\begin{enumerate}
  \item Untuk $\rho > 1$, misalkan $A_\rho$ luas (yakni \omterm{def:b3:measure:lebesgueouter}{ukuran
        Lebesgue}) himpunan \omterm{def:b3:topology:compact}{kompak} $K_\rho = \C\setminus
        g(\{\abs w > \rho\})$, yakni daerah yang dilingkupi
        kurva Jordan mulus $g(C_\rho)$. Dengan memakai rumus
        luas dari teorema Green--Riemann pada jilid Tahun ke-2
        --- yakni luas terlingkupinya adalah $\frac1{2\iu}
        \oint\bar\zeta\,\dd\zeta$ sepanjang batas yang
        berarah positif --- tunjukkan bahwa
        \[
        A_\rho = \frac{1}{2\iu}\int_{C_\rho}
        \overline{g(w)}\,g'(w)\,\dd w
        = \pi\Bigl(\rho^2 -
        \sum_{n\geq1}n\,\abs{b_n}^2\rho^{-2n}\Bigr) :
        \]
        substitusikanlah \omterm{thm:b3:residues:laurent}{deret Laurent} $\bar g$ dan $g'$ pada
        $C_\rho$ lalu integralkan suku demi suku (lewat kekonvergenan
        normalnya; sebab hanya hasil kali berfrekuensi nol yang
        bertahan).
  \item Lalu biarkan $\rho \to 1^+$ dan simpulkan \emph{teorema
        luasnya}:
        \[
        \sum_{n\geq1}n\,\abs{b_n}^2 \;\leq\; 1 .
        \]
        Khususnya $\abs{b_1} \leq 1$. Kapankah $\abs{b_1}
        = 1$?
\end{enumerate}

\medskip
\noindent\textbf{Bagian II --- Ketaksamaan $\abs{a_2} \leq
2$ milik Bieberbach.} Misalkan $f = z + a_2z^2 + \cdots \in \mathcal S$.
\begin{enumerate}[resume]
  \item Tunjukkan bahwa $f(z^2)/z^2$ bersifat \omterm{def:b3:holomorphic:holo}{holomorfik} dan bebas akar
        pada $\mathbb D$, dan memiliki akar kuadrat \omterm{def:b3:holomorphic:holo}{holomorfik}
        $\varphi$ dengan $\varphi(0) = 1$; lalu tetapkan $h(z) =
        z\varphi(z^2)$, sehingga $h(z)^2 = f(z^2)$. Lalu tunjukkan bahwa
        $h$ merupakan fungsi univalen yang \emph{gasal} pada $\mathbb
        D$ dengan uraian $h(z) = z + \frac{a_2}2z^3 +
        \cdots$. \emph{(Untuk keinjektifannya: $h(z)^2 = h(z')^2$
        memaksa $z^2 = z'^2$; lalu pakailah kegasalannya untuk menuntaskannya.)}
  \item Terapkan teorema luasnya pada $g(w) = 1/h(1/w) = w -
        \frac{a_2}{2w} + \cdots$ di $\{\abs w > 1\}$
        (periksalah keunivalenannya dan uraiannya), lalu simpulkan
        $\abs{a_2} \leq 2$.
  \item Tunjukkan bahwa \emph{fungsi Koebe}
        \[
        k(z) = \frac{z}{(1 - z)^2} = \sum_{n\geq1}n\,z^n
        \]
        termasuk $\mathcal S$, mempunyai $a_2 = 2$, dan memetakan
        $\mathbb D$ ke $\C\setminus\intoc{-\infty}{-\frac14}$
        \emph{(tulislah $k = \frac14\bigl[\bigl(\frac{1+z}{1 -
        z}\bigr)^2 - 1\bigr]$ lalu lacaklah petanya)}: sehingga semua
        ketaksamaan yang akan datang bersifat tajam.
\end{enumerate}

\medskip
\noindent\textbf{Bagian III --- Teorema seperempatnya.}
\begin{enumerate}[resume]
  \item Misalkan $f \in \mathcal S$ dan $c \notin f(\mathbb D)$.
        Tunjukkan bahwa
        \[
        F(z) = \frac{c\,f(z)}{c - f(z)}
        \]
        termasuk $\mathcal S$, lalu hitunglah koefisien
        keduanya: $A_2 = a_2 + \frac1c$.
  \item Terapkan Bieberbach pada $f$ maupun $F$: lalu simpulkan
        $\abs{\frac1c} \leq \abs{A_2} + \abs{a_2} \leq 4$,
        yakni $\abs c \geq \frac14$. \emph{Jadi setiap nilai yang dilewatkan
        bermodulus $\geq \frac14$}: sehingga $f(\mathbb D) \supseteq
        D(0, \frac14)$ --- yakni teorema seperempat Koebe. Lalu periksalah
        ketajamannya pada fungsi Koebenya.
  \item Turunkan sebuah taksiran kuantitatif pemetaan Riemannya: bahwa jika
        $\varphi \colon \Omega \to \mathbb D$ merupakan pemetaan Riemann
        \cref{thm:b3:conformal:rmt} di $z_0$, maka
        \[
        \frac{d\bigl(z_0, \partial\Omega\bigr)}{4}
        \;\leq\; \frac1{\varphi'(z_0)} \;\leq\;
        4\,d\bigl(z_0, \partial\Omega\bigr)
        \]
        --- lalu buktikanlah setidaknya ketaksamaan kirinya dengan menerapkan
        Koebe pada $\varphi^{-1}$ yang dinormalkan secara sesuai, dan
        yang kanannya lewat Schwarz yang diterapkan pada $\varphi$ pada cakram
        $D(z_0, d) \subseteq \Omega$.
\end{enumerate}

\medskip
\noindent\textbf{Bagian IV --- Perspektifnya.}
\begin{enumerate}[resume]
  \item Bieberbach menduga (1916) bahwa $\abs{a_n} \leq n$ untuk
        setiap $n$, dengan kesamaannya hanya bagi rotasi
        fungsi Koebenya; dan de Branges membuktikannya pada 1985. Periksalah
        dugaannya dengan tangan bagi fungsi Koebe dan
        rotasinya $\eu^{-\iu\theta}k(\eu^{\iu\theta}z)$. Lalu
        dorong uraian pertanyaan 4 satu suku lebih jauh: dengan menulis
        $h(z) = z + \frac{a_2}2z^3 + c_5z^5 + \cdots$, tunjukkan bahwa
        $c_5 = \frac{a_3}2 - \frac{a_2^2}8$ dan
        \[
        g(w) = w - \frac{a_2}{2}\,w^{-1} +
        \Bigl(\frac{3a_2^2}8 - \frac{a_3}2\Bigr)w^{-3} +
        \cdots ,
        \]
        sehingga teorema luasnya melahirkan ketaksamaan yang diperhalus
        $\bigl|\frac{a_2}2\bigr|^2 + 3\bigl|\frac{3a_2^2}8 -
        \frac{a_3}2\bigr|^2 \leq 1$. Lalu periksalah pada fungsi
        Koebenya (dengan $a_2 = 2$ dan $a_3 = 3$).
\end{enumerate}

\medskip
\noindent\textbf{Bagian V --- Teorema kesesatannya.}
Ketaksamaan Bieberbach, yang diangkut mengelilingi cakramnya lewat
automorfismanya, mengendalikan $f'$ di mana-mana. Tetapkan $f \in \mathcal
S$.
\begin{enumerate}[resume]
  \item (Transformasi Koebe) Untuk $z_0 \in \mathbb D$ misalkan
        $\varphi(z) = \frac{z + z_0}{1 + \bar z_0z}$, yakni sebuah
        automorfisma cakram (\cref{thm:b3:conformal:autdisc}) dengan
        $\varphi(0) = z_0$. Tunjukkan bahwa
        \[
        F(z) = \frac{f(\varphi(z)) - f(z_0)}
        {f'(z_0)\,\bigl(1 - \abs{z_0}^2\bigr)}
        \]
        termasuk $\mathcal S$ \emph{(sebab keunivalenannya
        diwarisi; lalu hitunglah $\varphi'(0) = 1 - \abs{z_0}^2$
        lalu periksa penormalannya; dan ingatlah $f' \neq 0$ bagi
        $f$ yang injektif,
        \cref{def:b3:conformal:conformal})}.
  \item Hitunglah koefisien kedua $A_2 = \frac12F''(0)$
        milik $F$:
        \[
        A_2 = \frac12\Bigl[\bigl(1 - \abs{z_0}^2\bigr)
        \frac{f''(z_0)}{f'(z_0)} - 2\bar z_0\Bigr] ,
        \]
        lalu turunkan dari Bieberbach (pertanyaan 4), untuk $z =
        r\eu^{\iu\theta}$, \emph{ketaksamaan
        fundamentalnya}:
        \[
        \Bigl|\,z\,\frac{f''(z)}{f'(z)} -
        \frac{2r^2}{1 - r^2}\Bigr|
        \;\leq\; \frac{4r}{1 - r^2} .
        \]
  \item Ambillah bagian realnya:
        \[
        \frac{2r^2 - 4r}{1 - r^2} \;\leq\;
        \operatorname{Re}\Bigl(z\,\frac{f''(z)}{f'(z)}\Bigr)
        \;\leq\; \frac{2r^2 + 4r}{1 - r^2} .
        \]
  \item Tunjukkan bahwa $\frac{\dd}{\dd t}\log\bigl|
        f'(t\eu^{\iu\theta})\bigr| = \frac1t
        \operatorname{Re}\bigl(z\frac{f''(z)}{f'(z)}\bigr)$
        di $z = t\eu^{\iu\theta}$ \emph{(sebab bagi fungsi
        $\mathcal C^1$ yang tak lenyap $g$ atas variabel real,
        $\frac{\dd}{\dd t}\log\abs g =
        \operatorname{Re}(g'/g)$)}, lalu integralkan batas
        pertanyaan 12 sepanjang sinarnya untuk memperoleh
        \emph{teorema kesesatannya}:
        \[
        \frac{1 - r}{(1 + r)^3} \;\leq\; \abs{f'(z)}
        \;\leq\; \frac{1 + r}{(1 - r)^3},
        \qquad \abs z = r .
        \]
  \item Turunkan \emph{teorema pertumbuhannya}:
        \[
        \frac{r}{(1 + r)^2} \;\leq\; \abs{f(z)} \;\leq\;
        \frac{r}{(1 - r)^2},
        \qquad \abs z = r
        \]
        \emph{(untuk batas atasnya: integralkan $f'$ pada ruas
        $[0, z]$; sedangkan untuk batas bawahnya: jika $\abs{f(z)} < \frac14$,
        maka ruas $[0, f(z)]$ terletak di $f(\mathbb D)$ menurut
        pertanyaan 7; lalu tariklah ia kembali lewat $f^{-1}$ ---
        yang \omterm{def:b3:holomorphic:holo}{holomorfik} menurut
        \cref{cor:b3:residues:openmapping} --- lalu batasi
        $\abs{f(z)} = \int_\gamma\abs{f'(\zeta)}
        \,\abs{\dd\zeta} \geq \int_0^r\frac{1 - t}{(1 +
        t)^3}\,\dd t$, dengan memakai $\abs{\dd\zeta} \geq
        \dd\abs\zeta$)}.
  \item Periksalah bahwa fungsi Koebenya mencapai kesamaan pada
        keempat batasnya, di $z = r$ untuk yang atas dan
        $z = -r$ untuk yang bawah: sehingga $k$ sekaligus
        merupakan anggota $\mathcal S$ yang paling memuai dan, di antipodenya, yang paling
        mengerut.
\end{enumerate}

\medskip
\noindent\textbf{Bagian VI --- Kekakuan ekstremalnya.} Pada setiap
ketaksamaan sejauh ini, kesamaannya mengenali fungsi Koebe sampai
rotasinya. Kita membuktikannya, lalu memanennya.
\begin{enumerate}[resume]
  \item Andaikan $f \in \mathcal S$ mempunyai $\abs{a_2} = 2$.
        Kejarlah kesamaannya melalui pertanyaan 2--4: sebab teorema
        luasnya memaksa $g(w) = w + b_0 + \eu^{\iu\alpha}/w$;
        lalu kegasalan $h$ membuat $g$ gasal, sehingga $b_0 = 0$; lalu balikkan
        untuk mencari $h$, kemudian $f$, lalu simpulkan bahwa
        \[
        f(z) = \eu^{-\iu\theta}k\bigl(\eu^{\iu\theta}z\bigr)
        \quad\text{dengan } \eu^{\iu\theta} =
        -\eu^{\iu\alpha} :
        \]
        sehingga rotasi fungsi Koebe merupakan satu-satunya
        anggota $\mathcal S$ yang $\abs{a_2} = 2$.
  \item Tunjukkan bahwa jika $f \in \mathcal S$ melewatkan sebuah nilai $c$
        yang $\abs c = \frac14$ persis, maka $f$ merupakan
        rotasi fungsi Koebe, lalu kenali $c =
        -\eu^{-\iu\theta}/4$ \emph{(lacaklah kesamaannya melalui
        rantai pertanyaan 7 $4 = \abs{1/c} = \abs{A_2 - a_2}
        \leq \abs{A_2} + \abs{a_2} \leq 4$)}: sehingga
        konstanta teorema seperempatnya tercapai hanya oleh keluarga
        ekstremalnya.
  \item (Koefisien secara murah) Gabungkan teorema
        pertumbuhannya dengan \omterm{thm:b3:holomorphic:analytic}{taksiran Cauchy}
        (\cref{thm:b3:holomorphic:analytic}) pada lingkaran
        $\abs z = 1 - \frac1n$ untuk membuktikan bahwa
        \[
        \abs{a_n} \;\leq\; \eu\,n^2 \qquad (n \geq 2) .
        \]
        (Menurut de Branges 1985: $\abs{a_n} \leq n$; sedangkan faktor
        $\eu n$-nya adalah harga perkakas dasarnya.)
  \item (Penyelimutan subcakramnya) Tunjukkan bahwa untuk setiap $0 < r <
        1$,
        \[
        f\bigl(D(0, r)\bigr) \supseteq
        D\Bigl(0, \frac{r}{(1 + r)^2}\Bigr),
        \]
        yang tajam bagi fungsi Koebenya, lalu pulihkan
        teorema seperempatnya saat $r \to 1^-$. \emph{(Sebab titik
        batas peta terbukanya $f(D(0,r))$ terletak pada
        $f(\partial D(0,r))$, sehingga bermodulus $\geq
        \frac{r}{(1+r)^2}$ menurut pertanyaan 14; dan sebuah ruas dari
        $0$ ke sebuah titik terlewat yang bermodulus lebih kecil harus
        menyeberangi batas itu.)}
  \item (Koebe di setiap titiknya) Misalkan $f$ univalen pada
        $\mathbb D$, yang tak harus ternormalkan, dan $z_0
        \in \mathbb D$. Buktikan bahwa
        \[
        \tfrac14\bigl(1 - \abs{z_0}^2\bigr)\abs{f'(z_0)}
        \;\leq\;
        d\bigl(f(z_0), \partial f(\mathbb D)\bigr)
        \;\leq\;
        \bigl(1 - \abs{z_0}^2\bigr)\abs{f'(z_0)}
        \]
        \emph{(untuk yang kiri: pakailah teorema seperempat yang diterapkan pada transformasi
        Koebe pertanyaan 10; sedangkan untuk yang kanan: Schwarz
        (\cref{thm:b3:conformal:schwarz}) yang diterapkan pada
        $\psi^{-1}\circ\hat g$, dengan $\hat g(w) =
        f^{-1}\bigl(f(z_0) + dw\bigr)$, $d$ jaraknya,
        dan $\psi$ sebuah automorfisma cakram yang mengirim $0$ ke
        $z_0$)}. Lalu mengapa $\partial f(\mathbb D)$ tak kosong?
  \item Periksalah pertanyaan 20 pada $f = k$ di $z_0 = r \in
        \intoo01$: hitunglah $d\bigl(k(r), \partial k(\mathbb
        D)\bigr) = \frac{(1+r)^2}{4(1-r)^2}$ lalu periksalah bahwa
        ketaksamaan kirinya berupa \emph{kesamaan}: sehingga fungsi
        Koebe menjenuhkan teoremanya sendiri di setiap titik
        $\intoo01$.
  \item (Moralnya) Dalam sepuluh baris: asas tunggal apa yang
        mendasari teorema luasnya, dan bagaimana Bieberbach,
        teorema seperempatnya, kesesatannya, pertumbuhannya, dan
        penyelimutannya semuanya mengalir darinya? Lalu bandingkan dengan
        dunia Schwarz--Pick pada
        \cref{thm:b3:conformal:schwarz}: sebab pada keduanya, satu
        ketaksamaan \omterm{def:b3:topology:interior}{interior} mengakukan seluruh geometrinya,
        dan ekstremalnya tunggal sampai rotasinya.

\end{enumerate}
\medskip
\noindent\textbf{Bagian VII --- Kekompakan, balikan, dan sebuah
uji kenyataan.}
\begin{enumerate}[resume]
  \item Tunjukkan bahwa kelas $\mathcal S$ bersifat \emph{\omterm{def:b3:topology:compact}{kompak}}
        terhadap kekonvergenan seragam lokalnya: sebab ia terbatas
        lokal menurut teorema pertumbuhannya, sehingga normal
        (\cref{thm:b3:conformal:montel}); dan limit seragam
        lokal anggota $\mathcal S$ kembali berada di
        $\mathcal S$ \emph{(sebab penormalannya beralih ke
        limitnya lewat kekonvergenan turunan Weierstrass;
        sedangkan keinjektifannya bertahan lewat Hurwitz,
        \cref{exo:b3:residues:8}, karena limitnya
        tak konstan)}. Lalu mengapa ini penting bagi soal
        ekstremal seperti soal Bieberbach?
  \item Untuk $f \in \mathcal S$, misalkan $g = f^{-1}$ yang terdefinisi
        di dekat $0$. Tunjukkan bahwa $g(w) = w - a_2w^2 + O(w^3)$, sehingga
        koefisien kedua balikannya menuruti batas tajam yang sama
        $\abs{A_2} = \abs{a_2} \leq 2$, dengan kesamaannya
        tepat bagi fungsi Koebe yang dirotasikan.
  \item Tentukan untuk $a \in \C$ yang mana polinomial $f(z)
        = z + az^2$ termasuk $\mathcal S$: tunjukkan bahwa $f$ bersifat
        injektif pada $\mathbb D$ jika dan hanya jika $\abs a \leq \frac12$
        \emph{(faktorkanlah $f(z_1) - f(z_2)$)}. Lalu simpulkan: bahwa untuk
        polinomial berderajat dua batas koefisien yang sebenarnya adalah
        $\frac12$, yakni empat kali lebih kecil daripada $2$ milik Bieberbach
        --- sehingga ekstremal $\mathcal S$ sungguh merupakan objek
        transenden, dan tak ada polinomial yang mendekatinya.

\end{enumerate}
\end{problem}
